The number of occurrences of highly cited papers is shown in Figure\ref{fig:gr3}.
Both the total number of highly cited papers and the number of newly appearing highly cited papers increase sharply after the cumulative generation spanning 1946-2010, which is thought to reflect the recent rapid rise in the number of indexed papers.

To assess the impact of highly cited papers, the proportion of highly cited papers among all papers is shown in Figure\ref{fig:gr4}.
Highly cited papers constitute only 0.0035\% of the total publications, yet they account for 1\% of citations, indicating a high concentration of citations. No consistent trend over time is observed, but a gentle fluctuation can be seen.

To examine the turnover of highly cited papers across cumulative generations, the proportion of newly highly cited papers within the set of highly cited papers is shown in Figure\ref{fig:gr5}. No consistent temporal trend is observed, and roughly a turnover of members ranging from 10\% to 60\% is evident.
